\def\versionname{学生练习版}
% !TEX program = xelatex
\documentclass[UTF8,11pt,a4paper,fontset=windows]{ctexart}
\usepackage{amsmath,amssymb,geometry,booktabs,array,fancyhdr,hyperref}
\geometry{left=20mm,right=20mm,top=19mm,bottom=19mm,headheight=14pt,footskip=10mm}
\hypersetup{unicode=true,hidelinks,pdfauthor={},pdftitle={直线和圆的方程 拔高讲义}}
\setlength{\parindent}{0pt}
\setlength{\parskip}{6pt}
\linespread{1.20}
\setlength{\emergencystretch}{2em}
\everymath{\displaystyle}
\raggedbottom
\pagestyle{fancy}
\fancyhf{}
\renewcommand{\headrulewidth}{0pt}
\fancyfoot[C]{\small 直线和圆的方程\quad\versionname\quad\thepage}
\ctexset{section={format=\Large\bfseries,number={},beforeskip=0pt,afterskip=10pt},subsection={format=\normalsize\bfseries,number={},beforeskip=9pt,afterskip=4pt}}
\newcommand{\qhead}[2]{\subsection{#1\texorpdfstring{\quad}{ }#2}}
\newcommand{\qsource}[1]{{\footnotesize 来源：#1。\par}}
\begin{document}

\section{2.1 直线的倾斜角与斜率}

\qhead{01}{线段遮挡与倾斜角区间}
\qsource{2025—2026学年度上学期七校协作体高二联考 数学}
取\(A\left(1,\allowbreak 2\right)\)、\(B\left(- 2,\allowbreak 0\right)\)、\(C\left(- 1,\allowbreak 4\right)\)。在经过\(C\)且不与线段\(AB\)相交的直线中，求所有存在的斜率\(k\)的范围。\par
思考题：求这些直线的倾斜角\(\alpha \)的范围，并说明竖直直线是否符合条件。\par
\vspace{12mm}\vfill

\qhead{02}{矩形中的动点斜率}
\qsource{2025—2026学年度上学期期中考试 高二数学}
矩形\(ABCD\)满足\(A\left(- 4,\allowbreak 4\right)\)、\(D\left(5,\allowbreak 7\right)\)。两条对角线的交点\(E\)位于第一象限，且\(E\)到\(y\)轴的距离为\(1\)。\(P\left(x,\allowbreak y\right)\)遍历边\(BC\)上使\(x\ne 0\)的点，求\(\frac{y}{x}\)的取值范围。\par
\vspace{12mm}\vfill

\mbox{}

\clearpage

\section{2.1 直线的倾斜角与斜率 续}

\qhead{03}{垂直条件与参数分类}
\qsource{牡丹江二中2025—2026学年度第一学期高二期中数学}
\(\ell \)的斜率为\(- \sqrt{3},\allowbreak \ell _{1}\)经过\(P\left(- 2,\allowbreak \sqrt{3}\right)\)和\(Q\left(m,\allowbreak 0\right)\)。若\(\ell _{1}\perp \ell ,\allowbreak m\)等于（　）。\par
\(A\)．\(- 2\)\qquad \(B\)．\(- 3\)\qquad \(C\)．\(- 4\)\qquad \(D\)．\(- 5\)\par
思考题：取消垂直条件。当\(m\)遍历实数时，分别写出\(\ell _{1}\)的倾斜角为锐角、直角、钝角的参数条件。\par
\vspace{12mm}\vfill

\mbox{}

\clearpage

\section{2.2 直线的方程}

\qhead{04}{平行条件中的重合陷阱}
\qsource{红桥高级中学2025—2026学年第一学期10月月考 高二数学}
\(\ell _{1}\)：\(\left(2a+ 1\right)x+ ay+ 1= 0,\allowbreak \ell _{2}\)：\(\left(a+ 2\right)x+ ay+ 2= 0\)。两条直线平行但不重合时，\(a\)的值为（　）。\par
\(A\)．\(- 1\)\qquad \(B\)．\(1\)或\(0\)\qquad \(C\)．\(0\)\qquad \(D\)．\(1\)\par
\vspace{12mm}\vfill

\qhead{05}{截距相等的完整讨论}
\qsource{牡丹江二中2025—2026学年度第一学期高二期中数学}
直线\(\ell \)经过\(A\left(1,\allowbreak 2\right)\)。分别求满足下列条件的全部直线方程。\par
（\(1\)）\(\ell \)与\(3x- 2y- 6= 0\)平行；（\(2\)）\(\ell \)在两个坐标轴上的截距相等。\par
\vspace{12mm}\vfill

\mbox{}

\clearpage

\section{2.2 直线的方程 续}

\qhead{06}{交点与截距倍数}
\qsource{2025—2026学年度上学期七校协作体高二联考 数学}
\(\ell \)通过\(x+ y- 1= 0\)与\(2x- y+ 4= 0\)的交点，且横截距是纵截距的\(2\)倍。写出\(\ell \)的全部可能方程。\par
\vspace{12mm}\vfill

\qhead{07}{被删掉的交点}
\qsource{2025—2026学年度上学期期中考试 高二数学}
令\(A= \left\{\left(x,\allowbreak y\right)\mid \frac{y- 3}{x- 2}= 2\right\},\allowbreak B= \left\{\left(x,\allowbreak y\right)\mid ax- y- 2= 0\right\}\)。使\(A\cap B= \varnothing \)的\(a\)可取（　）。\par
\(A\)．\(2\)\qquad \(B\)．\(- 2\)\qquad \(C\)．\(- \frac{5}{2}\)\qquad \(D\)．\(\frac{5}{2}\)\par
\vspace{12mm}\vfill

\mbox{}

\clearpage

\section{2.2 直线的方程 续}

\qhead{08}{中点条件确定截线}
\qsource{2024—2025学年上海市宝山区高境一中高二下学期3月月考数学}
直线\(\ell \)经过\(P\left(0,\allowbreak 1\right)\)，与\(\ell _{1}\)：\(2x+ y- 8= 0\)、\(\ell _{2}\)：\(x- 3y+ 10= 0\)分别交于\(B\)、\(A\)，且\(P\)是线段\(AB\)的中点。\par
（\(1\)）求\(\ell \)的方程；（\(2\)）设\(D\left(0,\allowbreak m\right)\)，且\(AD\parallel \ell _{1}\)，求\(\triangle ABD\)的面积。\par
\vspace{12mm}\vfill

\mbox{}

\clearpage

\section{2.3 直线的交点坐标与距离公式}

\qhead{09}{平行线距离与等距直线}
\qsource{2025年三新协同教研共同体高二联考 数学}
求\(2x- 3y+ 5= 0\)与\(4x- 6y+ 3= 0\)之间的距离。\par
思考题：写出到这两条直线距离相等的所有点组成的轨迹方程。\par
\vspace{12mm}\vfill

\qhead{10}{折线路程与取等点}
\qsource{牡丹江二中2025—2026学年度第一学期高二期中数学}
\(M\left(- 2,\allowbreak 3\right)\)、\(N\left(1,\allowbreak 1\right)\)固定，\(P\)在\(x\)轴上移动。\(\lvert PM\rvert + \lvert PN\rvert \)的最小值为（　）。\par
\(A\)．\(\sqrt{13}\)\qquad \(B\)．\(5\)\qquad \(C\)．\(\sqrt{5}\)\qquad \(D\)．\(\sqrt{17}\)\par
思考题：确定取得最小值时\(P\)的坐标。\par
\vspace{12mm}\vfill

\mbox{}

\clearpage

\section{2.3 直线的交点坐标与距离公式 续}

\qhead{11}{反射法叠加圆上动点}
\qsource{2025年三新协同教研共同体高二联考 数学}
\(P\)在直线\(y= 1\)上，\(Q\)在圆\(\left(x+ 1\right)^{2}+ \left(y+ 1\right)^{2}= 1\)上，\(A= \left(1,\allowbreak 0\right)\)。\(\lvert PA\rvert + \lvert PQ\rvert \)的最小值为（　）。\par
\(A\)．\(\sqrt{13}- 1\)\qquad \(B\)．\(\sqrt{13}+ 1\)\qquad \(C\)．\(2\sqrt{3}- 1\)\qquad \(D\)．\(2\sqrt{3}+ 1\)\par
\vspace{12mm}\vfill

\qhead{12}{直线束缺失位置与开端点}
\qsource{2025—2026学年度上学期期中考试 高二数学}
\(P\)遍历单位圆\(x^{2}+ y^{2}= 1,\allowbreak \lambda \)遍历实数。\(P\)到直线\(\left(2+ \lambda \right)x- \left(1+ \lambda \right)y- 6- 4\lambda = 0\)的距离可取范围为（　）。\par
\(A\)．\(\left[2\sqrt{2}- 1,\allowbreak 2\sqrt{2}+ 1\right]\)\qquad \(B\)．\(\left[0,\allowbreak 2\sqrt{2}+ 1\right]\)\par
\(C\)．\(\left[2\sqrt{2}- 1,\allowbreak 2\sqrt{2}+ 1\right)\)\qquad \(D\)．\(\left[0,\allowbreak 2\sqrt{2}+ 1\right)\)\par
\vspace{12mm}\vfill

\mbox{}

\clearpage

\section{2.4 圆的方程}

\qhead{13}{定比距离生成圆}
\qsource{牡丹江二中2025—2026学年度第一学期高二期中数学}
\(A= \left(- 2,\allowbreak 0\right)\)、\(B= \left(2,\allowbreak 0\right)\)，动点\(M\)满足\(\lvert MA\rvert = \sqrt{2}\lvert MB\rvert \)。下列结论哪些正确（　）？\par
\(A\)．\(M\)的轨迹为\(\left(x- 6\right)^{2}+ y^{2}= 32\)。\par
\(B\)．轨迹所围区域的面积为\(32\pi \)。\par
\(C\)．轨迹关于直线\(x- y- 6= 0\)对称。\par
\(D\)．点\(\left(7,\allowbreak 8\right)\)位于轨迹圆内部。\par
\vspace{12mm}\vfill

\qhead{14}{圆心定位与定长弦}
\qsource{2025年三新协同教研共同体高二联考 数学}
圆\(C\)经过\(\left(3,\allowbreak 1\right)\)、\(\left(1,\allowbreak 3\right)\)，圆心在\(x+ y- 6= 0\)上。\par
（\(1\)）求圆的方程；（\(2\)）过\(P\left(4,\allowbreak 0\right)\)的直线\(\ell \)截圆所得弦\(AB\)长为\(2\sqrt{3}\)，求全部\(\ell \)。\par
\vspace{12mm}\vfill

\mbox{}

\clearpage

\section{2.4 圆的方程 续}

\qhead{15}{切点定位与圆的轴对称}
\qsource{2025—2026学年度上学期期中考试 高二数学}
圆\(C\)经过\(A\left(2,\allowbreak - 1\right)\)，与\(x+ y= 1\)相切，圆心在\(y= - 2x\)上。过原点\(O\)的直线\(\ell \)截圆所得弦长为\(2\)。\par
（\(1\)）求所有\(\ell \)；（\(2\)）将圆\(C\)关于\(OA\)所在直线对称，求所得圆\(C'\)的方程。\par
\vspace{12mm}\vfill

\mbox{}

\clearpage

\section{2.4 圆的方程 续}

\qhead{16}{垂直直线束与轨迹完整性}
\qsource{2025—2026学年度上学期期中考试 高二数学}
\(m\in \mathbb{R} \)。直线\(x+ my= 0\)与\(mx- y- m+ 3= 0\)分别恒过\(A\)、\(B\)，交点为\(P\)。\(\lvert PA\rvert \cdot \lvert PB\rvert \)的最大值是（　）。\par
\(A\)．\(5\)\qquad \(B\)．\(10\)\qquad \(C\)．\(\frac{\sqrt{10}}{2}\)\qquad \(D\)．\(\sqrt{17}\)\par
思考题：求\(P\)的完整轨迹，写出必须剔除的点；再求取得上述最大值的全部\(m\)。\par
\vspace{12mm}\vfill

\qhead{17}{半圆约束下的分式范围}
\qsource{牡丹江二中2025—2026学年度第一学期高二期中数学}
实数\(x\)、\(y\)满足\(\sqrt{1- \left(y- 1\right)^{2}}= x- 1\)。则\(\frac{y+ 2}{x}\)的范围为（　）。\par
\(A\)．\(\left[\frac{4}{3},\allowbreak 4\right]\)\qquad \(B\)．\(\left[\frac{4}{3},\allowbreak 2\right]\)\qquad \(C\)．\(\left[\frac{4}{3},\allowbreak + \infty \right)\)\qquad \(D\)．\(\left[2,\allowbreak 4\right]\)\par
\vspace{12mm}\vfill

\mbox{}

\clearpage

\section{2.5 直线与圆 圆与圆的位置关系}

\qhead{18}{等距点个数的临界分类}
\qsource{2025年高考数学全国一卷}
在圆\(x^{2}+ \left(y+ 2\right)^{2}= r^{2}\)（\(r> 0\)）上，恰有两个点到\(y= \sqrt{3}x+ 2\)的距离等于\(1\)。\(r\)的范围为（　）。\par
\(A\)．\(\left(0,\allowbreak 1\right)\)\qquad \(B\)．\(\left(1,\allowbreak 3\right)\)\qquad \(C\)．\(\left(3,\allowbreak + \infty \right)\)\qquad \(D\)．\(\left(0,\allowbreak + \infty \right)\)\par
思考题：按\(r\)的变化，列出符合距离条件的点数，并解释两个临界半径为何不取。\par
\vspace{12mm}\vfill

\qhead{19}{直线束的弦长最值与缺失方向}
\qsource{2024—2025学年浙江省温州市高二上学期期末数学A卷}
圆\(C\)：\(\left(x- 1\right)^{2}+ \left(y- 1\right)^{2}= 16\)，直线\(\ell \)：\(\left(2m- 1\right)x+ \left(m- 1\right)y- 3m+ 1= 0,\allowbreak m\in \mathbb{R} \)。下列结论正确的是（　）。\par
\(A\)．\(\ell \)恒过定点\(\left(2,\allowbreak 1\right)\)\par
\(B\)．圆\(C\)被\(y\)轴截得的弦长为\(2\sqrt{15}\)\par
\(C\)．\(\ell \)截圆所得弦长存在最大值，此时\(\ell \)：\(2x+ y- 3= 0\)\par
\(D\)．\(\ell \)截圆所得弦长存在最小值，此时\(\ell \)：\(x- 2y- 4= 0\)\par
思考题：写出\(\ell \)截圆所得弦长的完整取值范围，说明两个端点是否能取到。\par
\vspace{12mm}\vfill

\mbox{}

\clearpage

\section{2.5 直线与圆 圆与圆的位置关系 续}

\qhead{20}{最小圆心角与最短弦}
\qsource{2026届华附 省实 广雅 深中高三四校期末联考 数学}
\(m\in \mathbb{R} \)，直线\(mx- y- 2m+ 3= 0\)与圆\(C\)：\(\left(x- 4\right)^{2}+ \left(y- 5\right)^{2}= 12\)相交于\(A\)、\(B\)。当\(\angle ACB\)最小时，求该角的余弦值。\par
思考题：求这时的\(m\)及弦长\(\lvert AB\rvert \)。\par
\vspace{12mm}\vfill

\qhead{21}{半圆的两个交点}
\qsource{2025—2026学年度上学期七校协作体高二联考 数学}
直线\(y= kx+ 3- k\)与上半圆\(y= \sqrt{1- x^{2}}\)有两个不同交点。\(k\)的范围是（　）。\par
\(A\)．\(\left(\frac{3}{4},\allowbreak + \infty \right)\)\qquad \(B\)．\(\left(\frac{3}{4},\allowbreak \frac{3}{2}\right]\)\par
\(C\)．\(\left(\frac{4}{3},\allowbreak \frac{3}{2}\right)\)\qquad \(D\)．\(\left(\frac{4}{3},\allowbreak \frac{3}{2}\right]\)\par
\vspace{12mm}\vfill

\mbox{}

\clearpage

\section{2.5 直线与圆 圆与圆的位置关系 续}

\qhead{22}{两圆公切线与相交弦}
\qsource{2025—2026学年度上学期七校协作体高二联考 数学}
圆\(C\)：\(x^{2}+ \left(y+ 2\right)^{2}= 4\)与圆\(D\)：\(\left(x+ 5\right)^{2}+ \left(y+ 2\right)^{2}= 9\)的公切线有（　）条。\par
\(A\)．\(1\)\qquad \(B\)．\(2\)\qquad \(C\)．\(3\)\qquad \(D\)．\(4\)\par
思考题一：求上述全部公切线方程。\par
思考题二：把圆\(D\)的半径改成\(\rho > 0\)，圆心不变。完整分类两圆的位置关系；\(\rho = 4\)时求公共弦所在直线和弦长。\par
\vspace{12mm}\vfill

\mbox{}

\clearpage

\section{2.5 直线与圆 圆与圆的位置关系 续}

\qhead{23}{切点弦中点的轨迹与距离}
\qsource{2025年三新协同教研共同体高二联考 数学}
\(Q\)在直线\(\ell \)：\(x- y- 4= 0\)上移动，向圆\(x^{2}+ y^{2}= 4\)作两条切线，切点为\(E\)、\(F\)，弦\(EF\)的中点为\(P\)。\(P\)到\(\ell \)的距离范围为（　）。\par
\(A\)．\(\left[\sqrt{3},\allowbreak 2\sqrt{3}\right]\)\qquad \(B\)．\(\left[\frac{\sqrt{3}}{2},\allowbreak \sqrt{3}\right)\)\par
\(C\)．\(\left[\sqrt{2},\allowbreak 2\sqrt{2}\right)\)\qquad \(D\)．\(\left[\sqrt{2},\allowbreak 2\sqrt{2}\right]\)\par
思考题：用圆的切线性质证明\(\lvert OP\rvert \cdot \lvert OQ\rvert \)为定值，求\(P\)的完整轨迹，并说明距离范围的两个端点能否取得。\(O\)为原点。\par
\vspace{12mm}\vfill

\mbox{}
\end{document}
